% Source: https://artofproblemsolving.com/community/c3h1521332_confusing_problem
% Stable identity: aops-topic-1521332-q01
For $n=p_1^{a_1}\cdots p_r^{a_r}$ the number of factors is $(a_1+1)\cdots(a_r+1)$. We show the maximum for $10\le n\le99$ is $12$.

\textbf{Four or more primes.} Impossible, since $2\cdot3\cdot5\cdot7=210>99$.

\textbf{Three primes} $p^aq^br^c$. More than $12$ factors would need $(a+1)(b+1)(c+1)\ge16$, i.e.\ exponents at least $3,1,1$ or $2,2,1$ in some order; the smallest such numbers are $2^3\cdot3\cdot5=120$ and $2^2\cdot3^2\cdot5=180$, too big. Exactly $12$ factors needs exponents $2,1,1$: $2^2\cdot3\cdot5=60$, $2^2\cdot3\cdot7=84$, $3^2\cdot2\cdot5=90$ (all others exceed $99$).

\textbf{Two primes} $p^aq^b$. At least $12$ factors needs exponents $(5,1)$, $(3,2)$ or more: $2^5\cdot3=96$ and $2^3\cdot3^2=72$ have $12$ factors, while $2^6\cdot3=192$, $2^4\cdot3^2=144$, $2^3\cdot3^3=216$, $2^2\cdot3^3=108$ and $3^5\cdot2$ are too big.

\textbf{One prime.} $2^6=64$ has only $7$ factors.

So the greatest number of factors is $12$, attained by $60,72,84,90,96$. Their sum is
\[
60+72+84+90+96=\boxed{402}.
\]
